\subsection{Discrete phase operator and transport grammar}

Let
\[
\omega=e^{2\pi i/3},
\qquad
U_6=\mathrm{diag}(1,\omega,\omega^2,-1,-\omega,-\omega^2).
\]
This makes the phase grammar concrete rather than rhetorical. The six admissible phases are shown as a finite hexagonal transport geometry in \Cref{fig:z6-circle}: the cusp values live on the same real axis as the positive-real phase, the admissible gap is visible directly, and the holonomy phase can be read as a triangle on the same discrete transport base.

\begin{remark}[Hexagonal spectral geometry]
The operator $U_6$ is a unitary circulant with eigenphases $e^{in\pi/3}$ for $n=0,\dots,5$ and is unitarily equivalent to the one-step translation operator on a six-site ring. The transport kernel $D_y=y\,\mathbf 1-\delta_{\mathrm{ph}}U_6$ is therefore the resolvent denominator of a finite quantum walk, equivalently a tight-binding transfer problem, on the discrete $C_6$ cycle. The regularized inverse introduced below plays the role of the corresponding Green kernel. CKM and PMNS mixings, mass hierarchies, and the Möbius quotients of \Cref{eq:mobius-quotient} are thus not loose analytic fits but spectral-theory statements on a regular hexagon. The near-pole hierarchy at $y\approx\delta_{\mathrm{ph}}$ becomes the resonance of a hexagonal resolvent rather than tuned parameter choice.
\end{remark}

\begin{figure}[t]
\centering
\scriptsize
\begin{tikzpicture}[scale=1.1,>=Latex]
\coordinate (O) at (0,0);
\foreach \name/\ang in {p0/0,p1/60,p2/120,p3/180,p4/240,p5/300} {
    \coordinate (\name) at ({2*cos(\ang)},{2*sin(\ang)});
}
\draw[thick, black!70] (p0) -- (p1) -- (p2) -- (p3) -- (p4) -- (p5) -- cycle;
\draw[gray!25, dashed] (O) -- (p0);
\draw[gray!25, dashed] (O) -- (p1);
\draw[gray!25, dashed] (O) -- (p2);
\draw[gray!25, dashed] (O) -- (p3);
\draw[gray!25, dashed] (O) -- (p4);
\draw[gray!25, dashed] (O) -- (p5);
\foreach \P in {p0,p1,p2,p3,p4,p5} {
    \filldraw[blue!60!black] (\P) circle (2pt);
}
\node[anchor=west, font=\scriptsize] at (2.18,0.12) {$1$};
\node[anchor=south west, font=\scriptsize] at (1.1,1.85) {$e^{i\pi/3}$};
\node[anchor=south east, font=\scriptsize] at (-1.1,1.85) {$e^{2i\pi/3}$};
\node[anchor=east, font=\scriptsize] at (-2.18,0.12) {$-1$};
\node[anchor=north east, font=\scriptsize] at (-1.1,-1.85) {$e^{-2i\pi/3}$};
\node[anchor=north west, font=\scriptsize] at (1.1,-1.85) {$e^{-i\pi/3}$};

\draw[->, thick, red!65!black] (1.15,1.05) arc[start angle=42,end angle=-42,radius=1.2];
\node[red!65!black, font=\scriptsize, anchor=west] at (2.55,0.25) {CP: $\theta\mapsto-\theta$};

\filldraw[orange!80!black] (0.67,0) circle (2pt);
\filldraw[orange!80!black] (1.33,0) circle (2pt);
\filldraw[green!60!black] (1.1998,0) circle (2.4pt);
\node[below, font=\scriptsize] at (0.67,-0.08) {$\frac13$};
\node[below, font=\scriptsize] at (1.33,-0.08) {$\frac23$};
\node[above, font=\scriptsize] at (1.1998,0.08) {$\delta_{\mathrm{ph}}$};
\draw[<->, thick, black!65] (1.1998,-0.32) -- (1.33,-0.32);
\node[below, font=\tiny] at (1.265,-0.34) {$\epsilon_f$};
\node[align=left, anchor=west, font=\scriptsize, text=black!70] at (2.55,-0.7)
    {positive-real branch\\near-pole lock at $y\approx\delta_{\mathrm{ph}}$\\gap $\epsilon_f=\frac23-\delta_{\mathrm{ph}}$};
\end{tikzpicture}
\caption{Finite transport geometry on the hexagon $C_6$. The six $U_6$ phases form the discrete
transport base, while the admissible cusps $y\in\{1,\frac13,\frac23\}$ and the branch pole
$\delta_{\mathrm{ph}}$ lie on the positive-real axis with gap
$\epsilon_f=\frac23-\delta_{\mathrm{ph}}$.}
\label{fig:z6-circle}
\end{figure}

The family-side holonomy phase is separated in \Cref{fig:holonomy-triangle} so that the transport
base and the flavor readout remain visually distinct.

\begin{figure}[t]
\centering
\scriptsize
\begin{tikzpicture}[>=Latex]
    \node[draw=black!65, fill=gray!10, rounded corners, text width=4.8cm, minimum height=0.85cm, align=center, font=\bfseries] at (0,1.15)
        {family-side holonomy on the transport grammar};
    \coordinate (A) at (0,0.1);
    \coordinate (B) at (-1.45,-1.8);
    \coordinate (C) at (1.45,-1.8);
    \draw[thick, orange!75!black] (A) -- (B) -- (C) -- cycle;
    \filldraw[orange!75!black] (A) circle (2pt) (B) circle (2pt) (C) circle (2pt);
    \node[above, font=\scriptsize] at (A) {$1$};
    \node[left, font=\scriptsize] at (B) {$2$};
    \node[right, font=\scriptsize] at (C) {$3$};
    \draw[->, thick, blue!60!black] (-0.2,-0.45) arc[start angle=110,end angle=15,radius=0.95];
    \node[draw=green!50!black, fill=green!8, rounded corners, text width=3.7cm, minimum height=0.95cm, align=center] at (0,-3.15)
        {$\delta_{\mathrm{CKM}}=\arg\tr\Pi_\triangle(\nabla_F^\star)$\\oriented holonomy triangle};
\end{tikzpicture}
\caption{Holonomy triangle on the same discrete transport grammar. The CKM phase is read as the
oriented three-cycle phase $\delta_{\mathrm{CKM}}=\arg\tr\Pi_\triangle(\nabla_F^\star)$ rather
than as an additional transport parameter.}
\label{fig:holonomy-triangle}
\end{figure}

\subsection{Regularized Yukawa kernels and path-length masses}

\[
P_{\mathrm{str}}:=P_{\Seam,+}P_{\mathrm{sing}}P_\Theta,
\qquad
P_{\Seam,+}:=\frac{\mathbf 1+\iC}{2},
\qquad
P_\Theta:=\frac12(\mathbf 1+Q_\Theta),
\qquad
P_{\mathrm{sing}}:=\frac13(\mathbf 1+Z_c+Z_c^2).
\]

\begin{definition}[Structural transport kernel]
\label{def:structural-transport-kernel}
For each fermion sector define the structural positive transport kernel by
\[
K_f^{\mathrm{str}}:=\bigl(P_{\mathrm{str}} D_f^\dagger D_f P_{\mathrm{str}}\bigr)^{-1/2},
\]
whenever the inverse exists on $\mathrm{Ran}(P_{\mathrm{str}})$. The final admissible kernel is
the restriction of this operator to the gap-stable branch selected below. The transport pole
$\delta_{\mathrm{ph}}$ is fixed by the structural stationarity of the exact determinant functional
\[
\mathcal F_{\mathrm{tr}}(\delta)
:=
\sum_{y\in\{1,1/3,2/3\}}
\Bigl[-\log\det\!\bigl(D_{y,\delta}^\dagger D_{y,\delta}+\varepsilon^2\bigr)\Bigr],
\]
with
\[
D_{y,\delta}:=y\,\mathbf 1-\delta U_6.
\]
On the closed branch one has
\begin{equation}
\label{eq:transport-bridge-offset}
\partial_\delta \mathcal F_{\mathrm{tr}}(\delta_{\mathrm{ph}})=0,
\qquad
\delta_{\mathrm{ph}}\in\left(\frac13,\frac23\right).
\end{equation}
\end{definition}

\begin{theorem}[Algebraic closure of the transport pole]
\label{prop:closed-form-transport-pole}
\label{thm:algebraic-transport-pole}
\label{thm:transport-pole-algebraic}
Let
\[
\mathcal F_{\mathrm{tr}}(\delta)
:=
\sum_{y\in\{1,1/3,2/3\}}
\Bigl[-\log\det\!\bigl(D_{y,\delta}^\dagger D_{y,\delta}+\varepsilon^2\bigr)\Bigr],
\qquad
D_{y,\delta}:=y\,\mathbf 1-\delta U_6.
\]
Write
\[
x:=\delta^2,
\qquad
\tau:=\varepsilon^2,
\qquad
z:=x^3=\delta^6.
\]
Then
\[
\det\!\bigl(D_{y,\delta}^\dagger D_{y,\delta}+\varepsilon^2\bigr)
=
Q_y(x,\tau),
\]
with
\[
Q_y(x,\tau)
=
\Bigl(x^2+2(\tau-y^2)x+(y^2+\tau)^2\Bigr)
\Bigl(x^2+(y^2+2\tau)x+(y^2+\tau)^2\Bigr)^2.
\]
Hence
\[
\partial_\delta \mathcal F_{\mathrm{tr}}(\delta)=0
\iff
\sum_{y\in\{1,1/3,2/3\}}
\frac{\partial_x Q_y(x,\tau)}{Q_y(x,\tau)}=0.
\]
On the closed zero-parameter branch one takes the intrinsic limit
\[
\tau=0
\]
after admissibility restriction. Then
\[
Q_y(x,0)=(y^6-z)^2,
\]
and the stationarity equation reduces to
\[
\frac{1}{1-z}
+
\frac{1}{1/729-z}
+
\frac{1}{64/729-z}
=0.
\]
Equivalently, for
\[
P(z):=
(z-1)\left(z-\frac{64}{729}\right)\left(z-\frac{1}{729}\right),
\]
one has
\[
P'(z)=0.
\]
The admissible transport pole is therefore the lower critical point of this cubic:
\[
z_\star
=
\frac{794-7\sqrt{9961}}{2187},
\qquad
\delta_{\mathrm{ph}}=z_\star^{1/6}
=
\left(\frac{794-7\sqrt{9961}}{2187}\right)^{1/6}
\approx \num{0.593277280133224}.
\]
As an expanded control form, one may equivalently write
\[
1594323\,\delta^{12}-1157652\,\delta^6+47449=0.
\]
\end{theorem}

\begin{proof}
In the Fourier basis of the six-site shift $U_6$, the operator
\[
D_{y,\delta}^\dagger D_{y,\delta}+\varepsilon^2
\]
has eigenvalues
\[
\lambda_{y,n}(x,\tau)
=
y^2+x-2y\sqrt{x}\cos\!\left(\frac{n\pi}{3}\right)+\tau,
\qquad n=0,\dots,5.
\]
Multiplying the $n=0,3$ factors and the conjugate pairs $(1,5)$ and $(2,4)$ gives exactly the
displayed factorization $Q_y(x,\tau)$. Since $x=\delta^2$ and $\delta\in(1/3,2/3)$ on the
admissible branch, one has $\delta>0$, so
\[
\partial_\delta\log Q_y(x,\tau)
=
2\delta\,\partial_x\log Q_y(x,\tau),
\]
and the stationarity condition is equivalent to the displayed $x$-equation. Setting $\tau=0$
collapses the factors to
\[
Q_y(x,0)
=(y^2-x)^2(y^4+xy^2+x^2)^2
=(y^6-z)^2.
\]
Therefore
\[
\sum_{y\in\{1,1/3,2/3\}}\frac{\partial_x Q_y(x,0)}{Q_y(x,0)}
=
-6x^2\sum_{y\in\{1,\frac13,\frac23\}}\frac{1}{y^6-z},
\]
so the stationary point condition reduces to the displayed rational equation. Let
\[
P(z):=(z-1)\left(z-\frac{64}{729}\right)\left(z-\frac{1}{729}\right).
\]
Since
\[
\frac{P'(z)}{P(z)}
=
\frac{1}{z-1}
+
\frac{1}{z-\frac{64}{729}}
+
\frac{1}{z-\frac{1}{729}},
\]
the stationary equation is equivalent to $P'(z)=0$. A real cubic with three ordered real roots has
exactly one critical point in each interval between consecutive roots, so the lower critical point
lies in
\[
z\in\left(\left(\frac13\right)^6,\left(\frac23\right)^6\right).
\]
Solving the quadratic $P'(z)=0$ gives
\[
z_\star=\frac{794-7\sqrt{9961}}{2187}.
\]
Hence
\[
\delta_{\mathrm{ph}}=z_\star^{1/6}\in\left(\frac13,\frac23\right).
\]
Expanding $P'(z)=0$ after substituting $z=\delta^6$ yields the displayed degree-$12$ control form.
\end{proof}

\begin{lemma}[Regularized transport pole stability]
\label{lem:regularized-transport-pole-stability}
Define the regularized stationarity function
\[
G(x,\tau):=
\sum_{y\in\{1,1/3,2/3\}}
\frac{\partial_x Q_y(x,\tau)}{Q_y(x,\tau)}.
\]
Let
\[
x_\star:=\delta_{\mathrm{ph}}^2,
\qquad
z_\star:=x_\star^3=\delta_{\mathrm{ph}}^6.
\]
Then
\[
\partial_xG(x_\star,0)\neq 0.
\]
Consequently there exists $\tau_0>0$ and a unique smooth branch
\[
x(\tau)\qquad (\tau\in[0,\tau_0))
\]
such that
\[
G(x(\tau),\tau)=0,
\qquad
x(0)=x_\star.
\]
Therefore
\[
\delta(\tau):=\sqrt{x(\tau)}=\delta_{\mathrm{ph}}+O(\tau)
\qquad (\tau\to 0).
\]
\end{lemma}

\begin{proof}
By the proof of \Cref{thm:transport-pole-algebraic},
\[
G(x,0)=6x^2\,\frac{P'(x^3)}{P(x^3)}.
\]
At $x=x_\star$ one has $P'(z_\star)=0$, so differentiating gives
\[
\partial_xG(x_\star,0)
=
18x_\star^4\,\frac{P''(z_\star)}{P(z_\star)}.
\]
Here $x_\star>0$, one has $P(z_\star)\neq 0$ because $z_\star$ is a critical point rather than a
root, and $P''(z_\star)\neq 0$ because the critical points of a cubic with three distinct roots are
simple. Hence $\partial_xG(x_\star,0)\neq 0$. The implicit-function theorem therefore yields a
unique smooth branch $x(\tau)$ through $(x_\star,0)$, and taking square roots gives the final
expansion for $\delta(\tau)$.
\end{proof}

\begin{definition}[Canonical cycle holonomy representation]
\label{def:canonical-cycle-holonomy}
On the rigid four-punctured family section
\[
X_f^\circ\cong \mathbb P^1\setminus\mu_4,
\qquad
\mu_4:=\{z\in\mathbb C:\ z^4=1\},
\]
fix a distinguished cycle basis
\[
(C_1,C_2,C_T)
\]
compatible with the lifted $D_4$ seam symmetry. Let
\[
F:=H^1(X_f^\circ,\mathbb C)\cong\mathcal H_F\cong\mathbb C^3.
\]
The canonical family holonomy datum is the determinant-normalized unitary representation
\[
\rho_F:\pi_1(X_f^\circ)\to SU(3)_F
\]
determined by the square symmetry of the puncture configuration, unimodularity, and the
admissible cusp assignment. Its concrete monodromy realization is forced by
\TFPTxref{thm:d4-monodromy-rigidity}.
\end{definition}
